% DRAFT — Jingyan Wei: check every item below, replace anything that is not
% actually yours, delete the ones you did not do, then commit this file under
% your own git account so that the history attributes it to you.
\clearpage
\begin{formbox}
\small
\textbf{Individual Contribution Statement --- Project 2 (PDE BVP)}\\[0.3em]
\textbf{Name:} Jingyan Wei \hfill \textbf{Student ID:} 2361411\\
\textbf{Role this project:} Variational formulation and the
$\star\star$ finite-element comparison.

\textbf{1. The specific things I did this project:}
\begin{enumerate}[leftmargin=1.5em,itemsep=1pt,topsep=2pt]
    \item I set up the $\mathcal P_1$ finite-element side of the comparison:
    the Friedrichs--Keller triangulation of the unit square, the element
    stiffness assembly from the constant element gradients, the congruent
    element-mass load vector for the source, and the strong Dirichlet
    elimination. This is the FE counterpart of the conservative FDM and it is
    what makes the two methods comparable in \S3.3.
    \item I fixed the matched-settings protocol of \S2.4 before any number was
    produced: the same number of unknowns, the same two norms evaluated on the
    same vertex set, the same relative residual tolerance with the same
    preconditioned CG, and the same manufactured solution. On the layered
    benchmark the two methods use different mesh parities by design, so I
    restricted the accuracy comparison to the smooth benchmark and compared
    structure and cost on the layered one.
    \item I produced the FE entries of the accuracy/cost tables and the
    corresponding curves: the FE nodal and $L^2$ errors over four meshes, the
    assembly and solve times, the non-zero count and the iteration count next
    to the FDM values, and the comparison against the direct sparse-LU solve.
    \item I prepared the interface-extension exploration archived in
    \path{project2/I-PINNs/}: two subnetworks with soft solution and flux
    constraints at the interface, run against the same layered benchmark. It is
    kept as documented material rather than as a claim of the report: the two
    independent networks are not a full reproduction of the published
    architecture, the layered exact solution depends on $x$ only, and the
    parameter count is roughly twice that of the single-network control.
\end{enumerate}

\textbf{2. The hardest conceptual obstacle I encountered:}
The Galerkin method enforces the weak equations on its own finite-dimensional
test space; it does not make the raw element flux pointwise continuous. On the
aligned layered benchmark the $\mathcal P_1$ solution is exact because the
piecewise-linear exact solution happens to belong to the space, which hides the
issue completely. Understanding that this exactness is a property of the
benchmark and not of the method was the obstacle, and it is why the report says
so explicitly instead of using the layered FEM run as evidence of flux
continuity.

\textbf{3. One piece of AI-generated output my team used, and what I changed:}
An AI-assisted draft of \S2.4 stated that a conforming discretisation
``automatically enforces interfacial flux continuity''. I corrected it: flux
continuity is a property of the continuous weak problem (it is the natural
transmission condition), while the discrete solution only satisfies the weak
balance on the FE space, and its element fluxes may jump. I rewrote the
paragraph, added the C\'ea/Aubin--Nitsche rate for $u\in H^2$ and the remark
that the aligned benchmark's exactness is exceptional (Appendix~A, Audit~3).

\textbf{4. Anything I should have asked for help with earlier:}
I should have asked for the two meshes to be fixed in a single configuration
file from the beginning; the FDM/FEM parity convention (odd for FDM, even for
FEM) changed twice and each change invalidated a set of tables.

\textbf{Signature:} \underline{\hspace{4.2cm}} \hfill
\textbf{Date:} \underline{\hspace{2.2cm}}
\end{formbox}
